\documentclass[11pt]{article}

\usepackage[margin=1in]{geometry}
\usepackage{amsmath, amssymb}
\usepackage{xcolor}
\usepackage{listings}
\usepackage{hyperref}
\usepackage{array}
\usepackage{booktabs}

\hypersetup{
    colorlinks=true,
    linkcolor=blue!60!black,
    urlcolor=blue!60!black,
    citecolor=blue!60!black
}

\lstdefinestyle{bugpython}{
  language=Python,
  basicstyle=\ttfamily\small,
  keywordstyle=\color{blue!60!black},
  stringstyle=\color{green!40!black},
  commentstyle=\color{gray!70!black},
  showstringspaces=false,
  breaklines=true,
  frame=single,
  rulecolor=\color{gray!30},
  columns=fullflexible,
  keepspaces=true
}

\title{SymPy Bug Report}
\author{}
\date{}

\begin{document}

\maketitle

\section{Bug 9: Finite Product Masks an Undefined Factor as Zero}

\subsection{Minimal Reproducer}

\medskip

\begin{lstlisting}[style=bugpython]
from sympy import pi, product, sin, symbols

n = symbols("n")
term = sin(pi*n)/(n - 1)

print("term =", term)
print("product(term, (n, 0, 2)) =", product(term, (n, 0, 2)))
print("term values =", [term.subs(n, i) for i in [0, 1, 2]])
\end{lstlisting}

\subsection{Observed Behavior}

\medskip

\begin{lstlisting}[style=bugpython]
term = sin(pi*n)/(n - 1)
product(term, (n, 0, 2)) = 0
term values = [0, nan, 0]
\end{lstlisting}

On SymPy 1.14.0, the finite product is evaluated as \(0\), even though direct
substitution of the individual factors shows that the middle factor is
\(\mathrm{nan}\). The zero endpoint factors do not make a product containing an
undefined factor well-defined.

\subsection{Expected Behavior}

The product should evaluate to \(\mathrm{nan}\), or otherwise remain unevaluated
with the singular factor exposed. It should not collapse to \(0\).

\subsection{Mathematical Explanation}

The product requested from SymPy is
\[
  \prod_{n=0}^{2} \frac{\sin(\pi n)}{n-1}.
\]
The three raw factors are
\[
  \frac{\sin(0)}{-1}=0,\qquad
  \frac{\sin(\pi)}{0}=\frac{0}{0},\qquad
  \frac{\sin(2\pi)}{1}=0.
\]
The middle factor is undefined. In SymPy's own arithmetic, direct multiplication
of the substituted factors gives \(\mathrm{nan}\), not \(0\). Therefore the
returned value is not an equivalent simplification, a harmless choice of
representation, or a convention for an empty product; it loses the domain
failure at \(n=1\).

\subsection{Source-Level Diagnosis}

The diagnosis locates the failure in \nolinkurl{sympy/concrete/products.py},
inside \texttt{Product.doit}, with the relevant code around line 256. The public
call \texttt{product} constructs a \texttt{Product} and calls
\texttt{Product.doit}. At the start of \texttt{doit}, definite product variables
are replaced by dummies with assumptions inherited from the concrete limits
before the product evaluator runs. For the limits \(0\) and \(2\), the helper
imported from \nolinkurl{sympy/concrete/summations.py} creates an
integer, nonnegative, real dummy.

\begin{lstlisting}[style=bugpython]
for xab in self.limits:
    d = _dummy_with_inherited_properties_concrete(xab)
    if d:
        reps[xab[0]] = d
...
did = self.xreplace(reps).doit(**hints)
\end{lstlisting}

After replacing \(n\) by this assumption-enriched dummy, \(\sin(\pi n)\)
simplifies as zero for an integer index. Consequently, the whole product
function becomes \(0\) before \texttt{\_eval\_product} sees the original
denominator \(n-1\) and the singular point \(n=1\). The diagnosis notes that the
downstream direct product evaluator in \nolinkurl{sympy/concrete/products.py}
would produce \(\mathrm{nan}\) for the original raw factors, but it is bypassed
because the pre-evaluated term is already zero. A plausible fix direction
identified by the diagnosis is to check original finite-range factors for
undefined values before accepting an identically zero product function, and more
generally to avoid allowing strengthened index assumptions to mask denominator
singularities in concrete products.

\subsection{Additional Failing Instances}

The independent verification covers a shifted family
\[
  \prod_{n=a-1}^{a+1}\frac{\sin(\pi n)}{n-a},
  \qquad a\in\{1,2,3,4,5,6\}.
\]
For every listed \(a\), the factor at \(n=a\) is \(0/0\), while the two adjacent
integer factors have zero numerator. The expected result is therefore
\(\mathrm{nan}\) across the family, because every product contains an undefined
factor.

\begin{center}
\renewcommand{\arraystretch}{1.3}
\begin{tabular}{@{}>{\raggedright\arraybackslash}p{0.44\linewidth}>{\raggedright\arraybackslash}p{0.23\linewidth}>{\raggedright\arraybackslash}p{0.23\linewidth}@{}}
\toprule
\textbf{Input / Expression} & \textbf{SymPy behavior} & \textbf{Expected behavior} \\
\midrule
\(\prod_{n=0}^{2}\sin(\pi n)/(n-1)\) & returns \(0\) & \(\mathrm{nan}\) \\
\(\prod_{n=1}^{3}\sin(\pi n)/(n-2)\) & returns \(0\) & \(\mathrm{nan}\) \\
\(\prod_{n=2}^{4}\sin(\pi n)/(n-3)\) & returns \(0\) & \(\mathrm{nan}\) \\
\(\prod_{n=3}^{5}\sin(\pi n)/(n-4)\) & returns \(0\) & \(\mathrm{nan}\) \\
\(\prod_{n=4}^{6}\sin(\pi n)/(n-5)\) & returns \(0\) & \(\mathrm{nan}\) \\
\(\prod_{n=5}^{7}\sin(\pi n)/(n-6)\) & returns \(0\) & \(\mathrm{nan}\) \\
\bottomrule
\end{tabular}
\end{center}

\subsection{Related Issues Assessment}

The deduplication check identified only general concrete-product documentation
as related background and did not find a public issue, pull request, Stack
Overflow question, mailing-list thread, or release-note entry for this exact
finite product or for the specific failure where integer-index assumptions
simplify \(\sin(\pi n)\) to zero before the denominator singularity is checked.
The verdict is a new failure mode with no public precedent. The bug remains
individually actionable because the reproducer is minimal, the affected call is
in SymPy's public concrete-product API, and the proposed regression test pins
the exact finite-range behavior that should be preserved.

\subsection{Proposed SymPy Regression Test}

\medskip

\begin{lstlisting}[style=bugpython]
import pytest
from sympy import Mul, nan, pi, product, sin, symbols


@pytest.mark.parametrize("a", [1, 2, 3, 4, 5, 6])
def test_finite_product_preserves_undefined_factor_before_zero_collapse(a):
    n = symbols("n")
    term = sin(pi*n)/(n - a)
    values = [term.subs(n, i) for i in range(a - 1, a + 2)]

    assert nan in values
    assert Mul(*values) is nan
    assert product(term, (n, a - 1, a + 1)) is nan
\end{lstlisting}

\end{document}
